% Partie foncier, extraite de ../foncier/note-foncier.tex puis retouchée
La France taxe beaucoup l'immobilier, et elle le taxe au mauvais endroit. Elle prélève lourdement au
moment où un logement change de mains : un ménage qui achète un logement ancien paie 5,8 à 6,3~\% du prix en
droits de mutation, en plus du prix, à l'instant précis où son apport est le plus contraint. Elle prélève
aussi beaucoup sur la détention, mais mal : la taxe foncière repose sur des valeurs locatives de 1970, qui
font payer proportionnellement plus les logements des communes pauvres que ceux des communes chères.

L'économie de ce choix est défavorable. Un impôt sur les ventes décourage les ventes : il retient dans leur
logement des ménages qui voudraient en changer, freine les déménagements vers un logement plus grand quand
la famille s'agrandit, et pèse d'abord sur ceux qui achètent, souvent jeunes. Un impôt récurrent sur la
valeur du sol, à l'inverse, ne décourage ni le travail ni la construction : le terrain existe de toute
façon. Il a un défaut, qui explique qu'on s'en serve si peu : il se voit, chaque année, sur un avis
d'imposition.

Cette note dresse l'état des lieux de tous les prélèvements assis sur l'immobilier, résume ce que la
recherche établit sur les droits de mutation et sur la taxe foncière, puis simule une bascule : diviser
par deux les droits de mutation, et reporter la recette sur une taxe foncière refondée, assise sur la
valeur de marché, avec un taux plus élevé sur le terrain que sur le bâti. L'objectif est une fiscalité
favorable à la croissance et à l'accession à la propriété, en particulier des jeunes ménages, sans
déstabiliser brutalement les budgets des propriétaires en place.

\subsection{État des lieux : 116 milliards, dont un quart sur les ventes}

Les prélèvements assis sur l'immobilier atteignent environ 116~Md€ en 2024, 4~\% du PIB
(figure~\ref{fig:inv}). La détention en fournit 69~Md€ : la taxe foncière sur les propriétés bâties
(42~Md€, 43,1 en 2025, dont environ 31,6~Md€ sur les logements), la taxe d'enlèvement des ordures ménagères
qui lui est adossée (9~Md€), la cotisation foncière des entreprises (8~Md€), l'impôt sur la fortune
immobilière (2,7~Md€) et les taxes d'habitation qui subsistent sur les résidences secondaires et les
logements vacants. Les transactions en fournissent 25~Md€ : droits de mutation (14,1~Md€ en 2024,
16,7 en 2025), TVA sur les logements neufs et les terrains à bâtir (7,1~Md€), impôt sur les plus-values
immobilières (2,8~Md€), contribution de sécurité immobilière (0,7~Md€). S'y ajoutent les droits de
succession et de donation sur l'immobilier (environ 12,5~Md€, estimation faute de ventilation
officielle), l'impôt sur les revenus locatifs (8,4~Md€) et la taxe d'aménagement (1,5~Md€).

\figbloc{fon_03_inventaire}{Les prélèvements assis sur l'immobilier en 2024}{fig:inv}

Selon l'OCDE, la France prélevait en 2023 en impôts périodiques sur l'immobilier 1,90~\% du PIB, le double
de la moyenne des pays membres (0,95~\%), et en impôts sur les transactions 0,72~\% contre 0,37~\%. Elle taxe
donc lourdement les deux. L'argument de cette note n'est pas que la détention serait peu taxée, mais qu'un
euro de droits de mutation coûte davantage à l'économie qu'un euro de taxe foncière, surtout s'il porte sur
le sol.

\textbf{Les droits de mutation.} Pour un logement ancien, ils cumulent la part départementale (4,50~\%),
la taxe communale (1,20~\%) et les frais d'assiette de l'État (2,37~\% de la part départementale), soit
5,81~\% du prix. La loi de finances pour 2025 a permis aux départements de porter leur part à 5~\% pour les
ventes conclues du 1\textsuperscript{er}~avril~2025 au 31~mars~2028 ; 88 collectivités sur une centaine
l'ont fait, ce qui porte le taux à 6,32~\%. Les primo-accédants achetant leur résidence principale sont
exclus de cette majoration ; la possibilité d'une réduction propre aux primo-accédants n'est utilisée que
par la Savoie.

\textbf{La taxe foncière.} Elle repose sur la moitié de la valeur locative cadastrale, établie selon les loyers
de 1970 et revalorisée depuis par des coefficients nationaux. La révision des valeurs des logements a été
reportée à 2031 par la loi de finances pour 2026. Rapportée à la valeur de marché, elle pèse 0,52~\% dans le
dixième des logements situés dans les communes les moins chères et 0,18~\% dans le dixième le plus cher
(figure~\ref{fig:tfv}) ; à Paris, 0,12~\%, malgré la hausse du taux de 13,5 à 20,5~\% en 2023. L'Insee aboutit au même constat par une autre voie : la taxe
foncière est régressive rapportée au patrimoine immobilier et progressive rapportée au revenu (André et
Meslin, 2023).

\subsection{Ce que dit la recherche}

\subsubsection*{Les droits de mutation réduisent les ventes et la mobilité}

Trois résultats sont solides. D'abord, les ventes réagissent fortement au calendrier : le mois qui précède
une hausse, elles bondissent de 26~\% en France en 2014 (Bérard et Trannoy, 2018), de 41~\% en Allemagne ;
ce sont des ventes avancées, pas des ventes créées. Ensuite, au-delà de ces effets de calendrier, la plupart
des études trouvent que chaque point de droits supprime environ 8~\% des ventes, dans une fourchette de 0 à
18~\% et sur des horizons de quelques mois à quatre ans (Dachis, Duranton et Turner, 2012, à Toronto ; Fritzsche et Vandrei, 2019, en
Allemagne ; Best et Kleven, 2018, au Royaume-Uni). Enfin, la mobilité des propriétaires baisse de 8 à 14~\%
par point (Pays-Bas, Finlande, Toronto). En Finlande, ce sont surtout les déménagements courts et les
déménagements vers plus grand qui disparaissent, pas les réductions de taille (Eerola et al., 2021) : la taxe
frappe le ménage qui veut une pièce de plus, ce qui est souvent le cas quand la famille s'agrandit (l'étude
n'observe pas les naissances).

Pour la France, les preuves sont plus minces et partagées. Après la hausse de 0,7~point de 2014, Bérard et
Trannoy trouvent une baisse des ventes de 6~\% sur trois mois, puis plus rien ; Bachelet et Poulhès (2017)
trouvent une baisse durable de 10 à 12~\%, de 20~\% dans les marchés détendus et nulle dans les marchés
tendus, et un prix hors droits inchangé : l'acheteur a supporté toute la hausse. En Allemagne, au contraire,
une hausse d'un point des droits réduit les prix d'environ 3~\% : les vendeurs la supportent (Dolls, Fuest,
Krolage et Neumeier, 2025). L'incidence d'une baisse est donc incertaine : à court terme, elle profite
plutôt aux acheteurs ; à plus long terme, elle passe en partie dans les prix au profit des vendeurs, et une
exonération ciblée sur les primo-accédants au Royaume-Uni a augmenté leurs achats de 11~\% tout en relevant
de 1~\% les prix qu'ils payaient (HMRC, 2023).

Le coût économique des droits de mutation croît vite avec leur taux, comme celui de tout impôt : autour de
1 à 2~\%, il représente 5 à 30~\% de la recette ; vers 4 à 5~\%, il approche 0,7 à 0,9~€ par euro
supplémentaire collecté (Buettner, en Allemagne ; Määttänen et Terviö, 2022). Au taux français, l'extrapolation
du modèle finlandais donne de l'ordre de 1 à 2~€ par euro, mais personne ne l'a estimé pour la France. Le seul
modèle français (Bontemps, Cherbonnier et Magnac, 2023) trouve un gain modeste à remplacer les droits par
un impôt foncier (+0,2~\% de consommation), qui s'inverse si l'on finance la baisse par l'impôt sur le
revenu : c'est la raison de la reporter sur la détention. Dans ce modèle, la baisse profite surtout aux
jeunes ménages aisés, dont l'accès à la propriété progresse, alors que celui des jeunes ménages modestes
recule.

\subsubsection*{La taxe foncière : un bon impôt mal assis}

Le principal inconvénient d'un impôt récurrent sur l'immobilier est la transition. Une hausse permanente se
capitalise dans les prix : le propriétaire présent au moment de l'annonce perd la valeur actuelle de toutes
les hausses futures, et les acheteurs suivants achètent moins cher (Koster et Pinchbeck, 2022, pour la
council tax anglaise). C'est un transfert des propriétaires en place vers les futurs acheteurs. Pour une
taxe assise sur le seul sol, l'ampleur de cette capitalisation reste discutée (Danemark : complète selon
Høj et al., 2018, nulle selon Nielsson et al., 2024). Un ménage âgé riche en patrimoine mais modeste en
revenu peut aussi avoir du mal à payer ; les reports de paiement facultatifs sont peu utilisés, sauf quand
ils sont proposés par défaut, comme au Danemark.

L'assiette est l'autre difficulté. Les évaluations de masse sont souvent régressives, parce qu'elles suivent
mal la valeur des quartiers chers (Berry, 2021 ; Avenancio-León et Howard, 2022). Philadelphie a corrigé
trente ans de gel par une réévaluation générale en 2014, accompagnée d'un abattement et d'un dispositif
pour les occupants de longue date. Pittsburgh, à l'inverse, a abandonné en 2001 son double taux après une
réévaluation brutale qui a fait bondir la valeur du sol de 81~\% sans lissage. Un plafond de hausse attaché
au propriétaire plutôt qu'au bien crée un verrou : en Californie, la Proposition~13 retient les ménages dans
leur logement (Ferreira, 2010).

Le double taux, plus élevé sur le terrain que sur le bâti, a un effet modeste mais positif sur la
construction : en Pennsylvanie, un rapport de quatre pour un a accru d'environ 3~points la croissance
décennale du nombre de pièces (Banzhaf et Lavery, 2010). La part de l'impôt qui porte sur le bâti décourage l'investissement là où l'offre
de logements peut réagir (Lutz, 2015). Le Bade-Wurtemberg est allé au bout de la logique : depuis 2025, sa
taxe foncière ne porte que sur la valeur du sol, établie par des valeurs de zone, avec un taux réduit de
30~\% pour le logement ; la plus haute juridiction fiscale allemande l'a jugée conforme en mai 2026, des recours
constitutionnels restant pendants. À Karlsruhe, à recette constante, les maisons individuelles paient
76~\% de plus et les copropriétés 20~\% de moins : l'assiette sol charge les grandes parcelles et allège la ville
dense.

\subsection{La proposition et ce qu'elle donne}

\begin{encadre}{Comment la bascule est chiffrée}
\textbf{Champ.} Les 30,5~millions de logements des particuliers (parc du recensement moins les HLM) de France
métropolitaine hors Alsace-Moselle, où les ventes ne sont pas publiées. Leur valeur est de 6~760~Md€ : prix
médian au m² des ventes 2023-2024 de la commune (sinon de l'intercommunalité ou du département), multiplié
par la surface médiane des logements vendus.

\textbf{Part du terrain.} Pour les maisons, prix au m² des terrains à bâtir multiplié par la surface de la
parcelle ; quand la commune compte trop peu de ventes de terrains, le prix départemental est ajusté selon
l'écart de prix des maisons (élasticité estimée de 1,54 sur 3~300 communes). Pour les appartements, valeur
du logement moins le coût de reconstruction du bâti, calé pour que l'ensemble retrouve la part du terrain
des comptes de patrimoine des ménages (45~\%). Elle va de 16~\% dans le dixième le moins cher à 67~\% dans le
plus cher.

\textbf{Charges actuelles.} Taxe foncière : rapport taxe/valeur de chaque commune selon l'Insee (2017),
vieilli selon la croissance de la taxe foncière de la commune de 2017 à 2024 (fichiers REI), appliqué aux
valeurs de 2024 et calé sur le total (23,6~Md€ pour ce champ). Droits de mutation : 5,81~\% des
ventes de logements anciens, 11,4~Md€ par an en moyenne en 2023-2024.

\textbf{Ce qui n'est pas modélisé.} Les effets sur les prix (capitalisation), l'âge et le revenu des
propriétaires, les locaux professionnels, les logements sociaux et les résidences secondaires en tant que
telles.
\end{encadre}

\textbf{Première étape : diviser par deux les droits de mutation.} La mesure porte sur les logements ; les
locaux professionnels, qui portent le total des droits à 16,7~Md€ en 2025, ne sont pas traités ici. Le taux
passe de 5,81~\% (taux de droit commun, une fois éteinte la majoration temporaire à 6,32~\% votée jusqu'en
mars 2028) à 2,9~\%. La perte de
recette est de 5,7~Md€ si les ventes ne réagissent pas, de 4,6~Md€ si chaque point de baisse augmente les ventes
de 6~\%, valeur centrale de la littérature (4,3 à 5,7~Md€ selon l'hypothèse). Une baisse de 30~\% coûterait 2,6
à 3,4~Md€. La baisse doit s'appliquer dès son annonce, sans délai : sinon les ventes s'arrêtent jusqu'à son
entrée en vigueur. La Finlande l'a fait en 2023 en la rendant rétroactive à la date du projet de loi. Par
prudence, on finance la baisse comme si les ventes ne réagissaient pas ; leur reprise, si elle a lieu, est un
gain.

\textbf{Deuxième étape : reporter la recette sur la taxe foncière.} Trois manières de le faire ont des
effets très différents (figures~\ref{fig:tfv} et~\ref{fig:var}).

La plus simple relève la taxe foncière actuelle de 24~\% partout (20~\% si les ventes réagissent) : entre 110
et 240~€ de plus par logement et par an. Elle ne bouleverse rien, mais elle aggrave l'injustice des
valeurs de 1970 : la taxe atteindrait 0,65~\% de la valeur dans les communes les moins chères contre 0,22~\%
dans les plus chères.

La plus ambitieuse remplace la taxe foncière par un impôt national à deux taux sur la valeur de marché :
0,73~\% sur le terrain et 0,18~\% sur le bâti. Elle supprime les distorsions de 1970, mais d'un coup : la taxe
augmenterait de 215~\% en moyenne dans le dixième des logements le plus cher et un logement sur quatre
verrait sa taxe augmenter de plus de moitié. Sans une transition de plusieurs décennies, c'est impossible.

\figbloc{fon_01_tf_valeur}{La taxe foncière rapportée à la valeur des logements, aujourd'hui et après la bascule}{fig:tfv}

\textbf{Notre proposition} combine les deux logiques. Chaque commune garde sa recette de taxe foncière et son
pouvoir de taux, mais l'assiette devient la valeur de marché, avec un taux sur le terrain égal à quatre fois
le taux sur le bâti, dans la fourchette pratiquée en Pennsylvanie. La perte de droits de mutation est
remplacée par un taux national unique de 0,19~\% sur la valeur des terrains, reversé aux départements selon
leurs pertes ; si les ventes repartent, il peut descendre vers 0,15~\%. Le principal effet économique est
de déplacer l'impôt du bâti vers le sol : la part de la taxe foncière qui porte sur les constructions passe
d'environ 14 à 7~Md€. Le taux national évite une difficulté
sérieuse : si chaque département remplaçait seul ses droits perdus, il lui faudrait un taux de 0,12~\% de la
valeur du sol à Paris mais de 0,85~\% dans la Creuse, où les terrains valent peu.

\figbloc{fon_02_variation}{Variation de la charge par logement et par an, selon le prix au m² de la commune}{fig:var}

La taxe foncière augmente de 3~\% dans le dixième des logements le moins cher et de 69~\% dans le plus
cher. En comptant la baisse des droits de mutation, rapportée à l'ensemble des logements de la commune, la
charge totale varie de $-$63 à +133~€ par logement et par an selon le décile. Rapportée à la valeur, la taxe
va de 0,54~\% dans les communes les moins chères à 0,30~\% dans les plus chères : l'écart actuel se resserre
sans être renversé. Les hausses se concentrent : un logement sur quatre (23~\%) voit sa taxe foncière
augmenter de plus de 300~€, et un sur vingt (5~\%, surtout à Paris et dans les communes littorales les plus
chères) la voit plus que doubler. Le partage entre maisons et appartements, lui, n'est pas robuste : il
dépend de la façon de mesurer le terrain. Avec notre méthode, les appartements paient davantage (+240~€
contre +150~€) ; si l'on mesure le terrain des maisons comme celui des appartements, par différence entre la
valeur et le coût de reconstruction, c'est l'inverse (+130~€ contre +230~€). Il faudra simuler l'assiette
légale, valeur de zone multipliée par la parcelle et répartie entre copropriétaires, avant de conclure.

Le gain d'efficacité est réel mais modeste. Si chaque euro de droits supprimé coûte en moyenne 0,3 à 0,9~€ à
l'économie, et chaque euro de taxe sur le sol presque rien, déplacer 5,7~Md€ rapporte de l'ordre de 1,5 à
4~Md€ par an, 0,05 à 0,15~\% du PIB : un ordre de grandeur, pas une estimation.

\figbloc{fon_04_primo}{Ce que gagne un ménage qui achète le logement médian de sa commune}{fig:primo}

\textbf{Pour les jeunes acheteurs, le gain dépend de qui supporte l'impôt} (figure~\ref{fig:primo}). Si
les prix ne bougent pas, comme après la hausse française de 2014, un ménage qui achète le logement médian de
sa commune paie entre 2~300~€ (un appartement à Saint-Étienne) et 22~700~€ (une maison à Arcachon) de droits
en moins : 8~100~€ pour un appartement à Lyon, 13~000~€ pour une maison à Nantes, 12~100~€ pour un
appartement à Paris. La hausse de taxe foncière qu'il paie ensuite (de $-$100 à +660~€ par an selon la
commune) ne rattrape cette économie qu'au bout de 17 à 23~ans dans les grandes villes, davantage ailleurs.
Le perdant est alors le propriétaire installé qui ne vend pas : il paie la hausse sans profiter de la baisse,
jusqu'à 660~€ par an pour l'appartement médian parisien, 22~000~€ en valeur actuelle.

Si au contraire les deux mouvements passent dans les prix, la baisse des droits les fait monter et la hausse
de taxe foncière les fait baisser ; à recette constante, le prix moyen bouge peu. La bascule transfère alors
de la valeur entre logements, des plus fonciers et des moins souvent vendus vers les autres, plutôt qu'entre
générations : l'appartement médian parisien perdrait environ 8~000~€, 2~\% de sa valeur. Dans les deux cas,
l'argument le plus solide en faveur des jeunes est la liquidité : les droits se paient comptant, sur
l'apport. Si l'apport représente 10 à 20~\% du prix, la baisse de 2,9~points en représente un septième à un
tiers.

\textbf{Pour les familles}, le gain tient au type de déménagement que les droits découragent le plus :
l'agrandissement. Avec des droits divisés par deux, passer d'un trois-pièces à un quatre-pièces à la
naissance d'un deuxième enfant coûte la moitié moins en impôt. Aucune étude
ne relie directement les droits de mutation à la natalité : l'effet attendu passe par le logement, pas par
une incitation directe.

\textbf{Pour les ménages les plus aisés}, l'effet net est limité. La taxe foncière augmente davantage dans
les communes chères (en moyenne 560~€ par an dans le dixième le plus cher), mais
les acheteurs de biens chers économisent davantage de droits. L'impôt sur la fortune immobilière n'est pas
modifié.

\subsection{Comment passer d'un système à l'autre}

\textbf{Construire l'assiette.} La France n'a pas de valeurs du sol par zone, comme les
\emph{Bodenrichtwerte} allemandes. Les ventes publiées (DVF) et leur version enrichie réservée aux acteurs
publics (DV3F) permettent de les établir par commune ou par quartier, en combinant les ventes de terrains et
les prix des logements, comme le fait cette note. L'administration fiscale établirait des valeurs de zone,
révisées tous les trois ans et lissées sur trois ans ; le propriétaire pourrait contester par expertise
quand sa valeur s'écarte nettement du marché, comme en Bade-Wurtemberg au-delà de 30~\%. Il faut écarter
la révision de la valeur à chaque vente, adoptée en commission le 9~octobre~2026 (amendement I-CF1872), puis
tombée avec le rejet de la première partie du projet de loi de finances : elle créerait le même verrou que la Proposition~13.

\textbf{Un calendrier.} Les valeurs du sol n'existent pas encore ; le Bade-Wurtemberg a mis cinq ans à les
établir (loi de 2020, premier impôt en 2025). Or la baisse des droits ne peut pas attendre : annoncée pour
plus tard, elle gèlerait les ventes jusqu'à son entrée en vigueur. Elle s'applique donc dès son vote et
elle est financée, le temps d'établir les valeurs, par une hausse uniforme et temporaire de la taxe foncière
actuelle (20 à 24~\%), remplacée ensuite par la nouvelle assiette. Celle-ci se substitue à la révision des
valeurs locatives reportée à 2031. Le taux national sur le terrain s'applique en totalité dès la première
année : il est faible (0,19~\% de la valeur du sol). Seule la nouvelle répartition de la taxe communale entre
logements est étalée, sur dix ans, hausses et baisses au même rythme, pour que la recette de chaque commune
reste constante chaque année. Les propriétaires âgés aux revenus modestes, déjà exonérés au-delà de 75~ans
sous condition de ressources, se verraient proposer par défaut un report de paiement portant intérêt,
remboursé à la vente ou à la succession. Le seuil de ressources, fondé sur le revenu fiscal de référence,
devra être redéfini si l'impôt sur le revenu et la CSG fusionnent.

\textbf{Garder l'autonomie locale.} Les communes votent leur taux sur la nouvelle assiette ; la loi fixe le
rapport entre le taux sur le terrain et le taux sur le bâti. Le Conseil constitutionnel n'exige pas de
compensation à l'euro près quand une recette locale est supprimée (décision 2019-796 DC) ; le taux national
sur le terrain, reversé aux départements, assure néanmoins cette compensation. Elle a un prix : les
départements échangent la moitié de leur principal impôt dont ils votent le taux contre une dotation, et
certains territoires paient plus qu'ils ne perdent (Paris environ 270~M€, les Hauts-de-Seine 130~M€). La taxe
d'enlèvement des ordures ménagères, assise sur la base de la taxe foncière, devra en être découplée.

\textbf{Limites.} Les valeurs sont des médianes communales, pas des valeurs de biens ; la part du terrain
des appartements est un calage sur les comptes nationaux, que l'Insee obtient lui-même par différence ; la
charge actuelle de chaque commune repose sur un rapport de 2017, vieilli par la croissance de sa taxe
foncière. L'effet de la taxe foncière française sur
les loyers et la construction n'a jamais été mesuré, et la répartition par âge de la bascule reste à
établir sur des données individuelles.
